\section{Convergent inverse charts at the critical endpoint}
\label{sec:charts}

We now resolve the critical singularity beyond its leading equivalent.
Throughout this section $u=\Uc(s)$, $s>2$, and the inverse is selected by
$x>0$ when $d>0$.  Its exact equation is
\eqref{eq:critical-normalform}.  All fractional powers are continued
from the positive axis in the sector supplied by \cref{lem:delta}.
A power of a function near $1$ uses the Taylor-series logarithm there.
Every chart below is holomorphic in finitely many independent small
coordinates; its convergent Taylor series becomes an exact generalized
power or logarithmic series after the indicated substitution.

\subsection{The range \texorpdfstring{$2<s<3$}{2<s<3}}

Put
\[
 \alpha=s-1,\qquad \beta=2-\alpha=3-s,\qquad
 C=\Gamma(1-s)>0,
\]
and use $B_s$ from \eqref{eq:K-noninteger}.  Thus
$K_s(x)=Cx^\alpha+x^2B_s(x)$.  Write
\[
 B_0=B_s(0)=\frac{1+\zeta(s-2)}2,\qquad
 B_1=B_s'(0)=-\frac{\zeta(s-3)}6.
\]
The two small exponents $\alpha-1$ and $\beta$ are both positive
and sum to one.  This relation removes the need for a third generator.

\begin{theorem}[A convergent chart with polynomial sectors]
\label{thm:chart-low}
Define
\begin{equation}\label{eq:chart-low-coordinates}
 q=\left(\frac{Ad}{C}\right)^{1/\alpha},\qquad
 p=q^{\alpha-1},\qquad z=q^\beta.
\end{equation}
There is a unique holomorphic germ $U(p,z)$ at $(0,0)$, with
$U(0,0)=1$, satisfying
\begin{equation}\label{eq:chart-low-equation}
 U^\alpha+\frac zC U^2 B_s(pzU)
 =1+\frac{pz}{A}U-\frac{C}{2A^2}p^2z.
\end{equation}
The critical inverse is exactly
\begin{equation}\label{eq:chart-low-inverse}
 x=qU(q^{\alpha-1},q^\beta).
\end{equation}
Moreover,
\begin{equation}\label{eq:chart-low-sectors}
 U(p,z)=1+\sum_{\ell\ge1}P_\ell(p)z^\ell,
 \qquad P_\ell\in\R[p],\qquad \deg P_\ell\le2\ell,
\end{equation}
with normal convergence on a sufficiently small bidisk.
\end{theorem}

\begin{proof}
Substitute $x=qU$ and $d=Cq^\alpha/A$ into
\eqref{eq:critical-normalform}, then divide by $Cq^\alpha$.
Since $q=pz$ and $q^\alpha=p^2z$, the result is
\eqref{eq:chart-low-equation}.  Its derivative in $U$ at
$(p,z,U)=(0,0,1)$ is $\alpha\ne0$.  The holomorphic implicit
function theorem gives the germ.  Substitution recovers the original
equation, and \cref{lem:delta} identifies the selected inverse.

At $z=0$ the solution is $U=1$.  In the perturbation side obtained
after isolating $U^\alpha-1$, every monomial $p^jz^\ell$ has
$\ell\ge1$ and $j\le2\ell$.  The three sources of such monomials
are $pz$, $p^2z$, and $p^jz^{j+1}$ from the analytic tail.
Induction on the $z$-degree in the implicit coefficient recursion
preserves this support condition.  It proves the polynomial assertion;
holomorphy proves convergence.
\end{proof}

There is a finite formula for every sector polynomial.  Put
\begin{align*}
 F_0(h)&=(1+h)^\alpha-1,\\
 \Psi(p,z,h)&=\frac{pz}{A}(1+h)-\frac{Cp^2z}{2A^2}
 -\frac zC(1+h)^2B_s(pz(1+h)).
\end{align*}
Then
\begin{equation}\label{eq:chart-low-finite}
 P_\ell(p)=\sum_{k=1}^{\ell}\frac1k
 [z^\ell h^{k-1}]
 \left\{\frac h{F_0(h)}\Psi(p,z,h)\right\}^{k}.
\end{equation}
Indeed, insert an auxiliary multiplier on $\Psi$ in $F_0(h)=\Psi$,
rewrite this as a Lagrange equation for $h$, and extract its coefficients.
Since $\Psi$ is divisible by $z$, only $k\le\ell$ can contribute.
This finite extraction agrees with the Taylor coefficients of the
already established analytic germ; no convergence assertion about the
auxiliary multiplier is needed.

For example,
\begin{align}
 P_1(p)&=-\frac{B_0}{\alpha C}+
          \frac p{\alpha A}-\frac{Cp^2}{2\alpha A^2},
          \label{eq:chart-low-P1}\\
 P_2(p)&=\frac1\alpha\left\{
 \left(\frac pA-\frac{2B_0}C\right)P_1(p)-\frac{B_1p}C
 -\frac{\alpha(\alpha-1)}2P_1(p)^2\right\}.
          \label{eq:chart-low-P2}
\end{align}
The whole ladder at $p=0$ has the closed form
\begin{equation}\label{eq:chart-low-pure}
 P_\ell(0)=
 \frac{(-B_0/C)^\ell}{\alpha\,\ell!}
 \frac{\Gamma((2\ell+1)/\alpha)}
 {\Gamma(1+(1+\ell\beta)/\alpha)}.
\end{equation}
To prove it, set $p=0$ in the chart and put $w=U^\alpha$.
The equation becomes $w+(B_0/C)zw^{2/\alpha}=1$;
ordinary Lagrange inversion applied to $w^{1/\alpha}$ gives
\eqref{eq:chart-low-pure}.

The terms of total Taylor degree at most two are
\begin{equation}\label{eq:chart-low-first}
 U=1-\frac{B_0}{\alpha C}z+\frac1{\alpha A}pz
 +\frac{(5-\alpha)B_0^2}{2\alpha^2C^2}z^2
 +O((|p|+|z|)^3).
\end{equation}
Their asymptotic order after substitution depends on $s$; equal powers
must be collected when the two exponents are rationally related.
In particular, if $B_0\ne0$,
\[
 x=q-\frac{B_0}{\alpha C}q^{1+\beta}+o(q^{1+\beta}).
\]
There is a genuine cancellation possibility.  If $\zeta(s-2)=-1$,
then $B_0=0$ and the entire ladder \eqref{eq:chart-low-pure} vanishes.
In this case $U(p,0)=U(0,z)=1$, so $U-1$ is divisible by $pz=q$.
Equation \eqref{eq:chart-low-P1} therefore gives instead
\begin{equation}\label{eq:chart-low-cancellation}
 x=q+\frac{q^2}{\alpha A}+o(q^2).
\end{equation}
No uniqueness claim about an order satisfying this cancellation condition
is needed for the result.

\subsection{An exact Lambert chart at order three}

Let $s=3$, $A=\zeta(2)$, and define the analytic tail $C_3$ by
\[
 K_3(x)=\frac{x^2}{2}\left(\log\frac1x+\frac52\right)+x^3C_3(x).
\]
By \eqref{eq:K3},
$C_3(x)=1/12-x/288+O(x^3)$; its Taylor series converges for
$|x|<2\pi$.

\begin{theorem}[Convergent marginal inverse]
\label{thm:chart-three}
For small positive $d$, put
\begin{equation}\label{eq:three-Lambert-core}
 \Lambda=-\frac12W_{-1}(-4Ad\ee^{-5}),\qquad
 \tau=\sqrt{\frac{2Ad}{\Lambda}}.
\end{equation}
Thus $\Lambda=\log(1/\tau)+5/2$ and $\tau^2\Lambda=2Ad$.
There is a holomorphic germ $U(p,w)$ with $U(0,0)=1$ such that
\begin{equation}\label{eq:three-chart-inverse}
 x=\tau U(\tau\Lambda,\Lambda^{-1}).
\end{equation}
Its defining equation is
\begin{equation}\label{eq:three-chart-equation}
 U^2(1-w\log U)+2pw^2U^3C_3(pwU)
 =1+\frac{pw}{A}U-\frac{p^2w}{4A^2}.
\end{equation}
The germ $U-1$ is divisible by $pw$, and
\begin{equation}\label{eq:three-chart-correction}
 x=\tau+\tau^2\frac{\Lambda/A-1/6}{2\Lambda-1}
       +O(\tau^3\Lambda).
\end{equation}
\end{theorem}

\begin{proof}
The branch $W_{-1}$ in \eqref{eq:three-Lambert-core}, with the
standard convention of \cite{DLMFLambert}, is selected
because $\Lambda\to+\infty$.  The core solves
$\tau^2\{\log(1/\tau)+5/2\}=2Ad$ exactly.
Put $p=\tau\Lambda$, $w=1/\Lambda$, so that
$\tau=pw$ and $d=p^2w/(2A)$.  Substitution of $x=\tau U$ in
\eqref{eq:critical-normalform}, followed by division by $Ad$,
gives \eqref{eq:three-chart-equation}.  Its derivative in $U$ at
the origin is $2$, so the holomorphic implicit function theorem applies.
At $p=0$ or $w=0$, uniqueness gives $U=1$, proving divisibility.

At fixed $w$, the coefficient of $p$ is
\[
 [p]U(p,w)=\frac{w/A-w^2/6}{2-w}.
\]
The remaining terms are $O(p^2w)$ uniformly for small $w$.
Multiplication by $\tau$ and physical substitution yield
\eqref{eq:three-chart-correction}.  The real inverse agrees by
\cref{lem:delta}, and its sectorial continuation fixes the complex branch.
\end{proof}

Here too every coefficient is computable by finite extraction.
For $U=1+\sum_{n\ge1}u_n(w)p^n$, set
\begin{align*}
 F_0(h,w)&=(1+h)^2\{1-w\log(1+h)\}-1,\\
 \Psi(p,w,h)&=\frac{pw}{A}(1+h)-\frac{p^2w}{4A^2}
 -2pw^2(1+h)^3C_3(pw(1+h)).
\end{align*}
Then
\begin{equation}\label{eq:three-chart-finite}
 u_n(w)=\sum_{k=1}^{n}\frac1k[p^nh^{k-1}]
 \left\{\frac h{F_0(h,w)}\Psi(p,w,h)\right\}^{k}.
\end{equation}
These are rational functions of $w$, with their only possible finite
pole at $w=2$, of order at most $2n-1$.  Indeed,
$F_0(h,w)/h=2-w+O(h)$; extracting $h^{k-1}$ from its negative
$k$th power requires at most $2k-1$ inverse factors of $2-w$.
The assertion concerns coefficient functions of a chart proved near
$w=0$; it does not assert a singularity of a global physical inverse.

\subsection{Noninteger and integer orders above three}

For noninteger $s>3$, put
\[
 \mu=s-3>0,\quad C=\Gamma(1-s),\quad
 B=\frac{1+\zeta(s-2)}2>0,\quad q=\sqrt{Ad/B}.
\]
Equation \eqref{eq:K-noninteger} reads
$K_s(x)=x^2B_s(x)+Cx^{\mu+2}$, with $B_s(0)=B$.

\begin{theorem}[Convergent power and logarithmic charts]
\label{thm:chart-high}
For noninteger $s>3$, the critical inverse is
\begin{equation}\label{eq:chart-high-inverse}
 x=qU(q,q^\mu),
\end{equation}
where $U$ is holomorphic at $(0,0)$, $U(0,0)=1$, and
\begin{equation}\label{eq:chart-high-equation}
 U^2\frac{B_s(qU)}B+\frac CB zU^{\mu+2}
 =1+\frac qA U-\frac{Bq^2}{2A^2}.
\end{equation}
For an integer $m\ge4$, let
\[
 \nu=m-3,\quad A=\zeta(m-1),\quad
 B=\frac{1+\zeta(m-2)}2,\quad
 c_m=\frac{(-1)^{m-1}}{(m-1)!}.
\]
Define the analytic function $B_m$ by
\[
 K_m(x)=x^2B_m(x)+c_mx^{\nu+2}\log(1/x),\qquad B_m(0)=B.
\]
With $q=\sqrt{Ad/B}$, the integer-order inverse is
\begin{equation}\label{eq:chart-integer-inverse}
 x=qU(q,q^\nu\log(1/q)),
\end{equation}
where $U(0,0)=1$ is holomorphic and satisfies
\begin{equation}\label{eq:chart-integer-equation}
 U^2\frac{B_m(qU)}B+\frac{c_m}{B}zU^{\nu+2}
 -\frac{c_m}{B}q^\nu U^{\nu+2}\log U
 =1+\frac qA U-\frac{Bq^2}{2A^2}.
\end{equation}
Both formulas are exact convergent two-generator expansions.
\end{theorem}

\begin{proof}
In each case substitute $x=qU$, $d=Bq^2/A$ into
\eqref{eq:critical-normalform} and divide by $Bq^2$.
For the integer case use
$\log(1/(qU))=\log(1/q)-\log U$.  The respective equations follow.
Their derivatives in $U$ at the origin are $2$, proving the
holomorphic charts.  The positive and sectorial branch identifications
are the same as in \cref{thm:chart-low}.
\end{proof}

Every mixed coefficient has a finite Lagrange formula of the same kind
as \eqref{eq:chart-low-finite}.  More precisely, let
$F_0(h)=(1+h)^2-1$ and move all remaining terms of the relevant chart
equation to its right side, denoted $\Psi(q,z,h)$.
Then
\begin{equation}\label{eq:chart-high-finite}
 [q^nz^\ell](U-1)=
 \sum_{k=1}^{n+\ell}\frac1k[q^nz^\ell h^{k-1}]
 \left\{\frac h{F_0(h)}\Psi(q,z,h)\right\}^{k}.
\end{equation}
The constant term of $\Psi$ in $(q,z)$ vanishes, so the sum is finite.
For the noninteger case the first bivariate terms are
\begin{equation}\label{eq:chart-high-first}
 U(q,z)=1-\frac C{2B}z+
 \left(\frac1{2A}+\frac{\zeta(s-3)}{12B}\right)q
 +O(|q|^2+|qz|+|z|^2),
\end{equation}
and the pure fractional sectors are
\begin{equation}\label{eq:chart-high-pure}
 [z^k]U(0,z)=\frac{(-C/B)^k}{2k!}
 \frac{\Gamma(((\mu+2)k+1)/2)}{\Gamma((\mu k+3)/2)}.
\end{equation}
The latter follows by Lagrange inversion of
$U^2+(C/B)zU^{\mu+2}=1$.

At integer order, collecting equal powers in the absolutely convergent
double series gives
\begin{equation}\label{eq:chart-integer-degree}
 x=q\sum_{j\ge0}q^jQ_j(\log(1/q)),\qquad
 \deg Q_j\le\lfloor j/\nu\rfloor.
\end{equation}
This bound is attained at every $j=\nu k$: the highest logarithmic
coefficient comes from the $q$-constant part of the chart and equals
\begin{equation}\label{eq:chart-integer-top}
 [L^k]Q_{\nu k}(L)=\frac{(-c_m/B)^k}{2k!}
 \frac{\Gamma(((\nu+2)k+1)/2)}{\Gamma((\nu k+3)/2)}\ne0.
\end{equation}
In particular, order four requires the resonant constant
$B_4'(0)=-H_3/6=-11/36$, and yields
\begin{equation}\label{eq:chart-four}
 x=q+\frac{q^2}{12B}\log\frac1q
 +\left(\frac1{2A}+\frac{11}{72B}\right)q^2
 +O(q^3\log^2(1/q)).
\end{equation}
The harmonic constant is essential: directly substituting $s=4$ in
the noninteger coefficient \eqref{eq:chart-high-first} would encounter
$\zeta(1)$ and lose the resonant cancellation.
